\documentclass[10pt,a4paper]{amsart}
\usepackage[margin=19mm]{geometry}
\usepackage{amsmath,amssymb,mathtools}
\usepackage[scr=rsfs,cal=euler]{mathalfa}
\usepackage{xfrac}
\usepackage[unicode,hidelinks,bookmarks=false]{hyperref}
\newcommand{\ie}{i.e.\ }
\newcommand{\cf}{cf.\ }
\newcommand{\resp}{respectively\ }
\newcommand{\Subsection}{\subsection}
\providecommand{\nobreakdash}{\nobreak\hbox{-}}
\setlength{\emergencystretch}{2em}
\multlinegap0pt
\usepackage{dsfont}
\usepackage{amssymb}
\usepackage{float}
\usepackage{xcolor}
\newtheorem{theorem}{Theorem}
\title{Hamiltonian paths in the permutation digraphs $P(n,n-2)$}
\author{Jiaxin Guo$^{a}$, Ming Duan$^{a}$, Jie Xue$^{b}$}
\date{Source: arXiv:2608.15287v1 (CC BY 4.0)}
\begin{document}
\maketitle

\noindent\emph{Source and licence.} Hamiltonian paths in the permutation digraphs $P(n,n-2)$. Version arXiv:2608.15287v1 (CC BY 4.0), distributed by arXiv under the Creative Commons Attribution 4.0 International licence, http://creativecommons.org/licenses/by/4.0/, as recorded in the arXiv licence metadata for this exact version and shown on the abstract page https://arxiv.org/abs/2608.15287v1. Full licence notice: \texttt{licenses/arxiv-2608.15287-CC-BY-4.0.txt}.\par\smallskip

\noindent\textit{Editorial scope.} Counted unit: the article's theorem thm:root-deleted (source0.tex lines 514--520). The article's complete original proof of it is delivered with it (source0.tex lines 522--599, one \texttt{proof} environment of 78 lines). No mathematics is written, repaired or re-derived: the statement and the proof are the article's own lines, in the article's own notation, and no retained line is substituted or re-flowed.  No further source-local context is needed: every cross-reference of every retained line resolves inside the two delivered spans.  Nothing was omitted from any retained span, so the package carries no omission record.  No retained line contains a citation, so the wrapper carries no bibliography and no bibliography entry.  No retained line contains a figure environment, an image-inclusion command or drawing code, so no image is delivered and the package declares no asset dependency.  Editorial formatting only, in addition: an AMS \texttt{amsart} preamble that restates the article's own declarations and macros used by the retained lines, namely the environments \texttt{theorem} on separate counters, together with the standard packages those lines need.  The retained environments print in this document as Theorem 1 in order of appearance, whereas the article prints the same items with the numbering of its own full document.  The complete native source of the article as distributed in the arXiv e-print for this exact version is bundled unchanged as \texttt{sources/207729/source0.tex}.
\begin{theorem}\label{thm:root-deleted}
Let $m\geq4$, let $T$ be a $2$-tree on $X$, let $\Omega$ be a parity class
of bijections $X\to Y$, and let $r\in\Omega$. Then $K(T,\Omega)-r$ contains
a spanning loose hypertree. Equivalently, $K(T,\Omega)$ contains a spanning
loose hyperforest whose components are the isolated vertex $r$ and a loose
hypertree on $\Omega\setminus\{r\}$.
\end{theorem}

\begin{proof}
We use induction on $m$. For $m=4$, relabeling $X$ and $Y$ and
left-composing all bijections with $r^{-1}$ reduces the problem to
$X=Y=[4]$, $\Omega=A_4$, $r=1234$ in one-line notation, and triangles
$\{1,2,3\}$ and $\{2,3,4\}$. The following five hyperedges belong to
$K(T,A_4)$:
\[
\begin{array}{ll}
 E_1=\{1342,3412,4132\},&
 E_2=\{1423,2143,4213\},\\
 E_3=\{2143,2314,2431\},&
 E_4=\{2431,3241,4321\},\\
 E_5=\{3124,3241,3412\}.&
\end{array}
\]
In the order $E_2,E_3,E_4,E_5,E_1$, each edge after the first meets the
preceding union in one vertex. Their union is $A_4\setminus\{1234\}$, so
they form the required loose hypertree.

Let $m\geq5$, and choose a simplicial vertex $x$ of $T$ such that $T-x$ is a
$2$-tree. Write $\{x,a,b\}$ for the unique triangle containing $x$ and take
its generator to be
$
 \tau=(x\ a\ b).
$
For $j\in Y$, define the layer
$
 L_j=\{\pi\in\Omega:\pi(x)=j\}.
$
Every triangle generator other than $\tau$ fixes $x$. After restriction to
$X\setminus\{x\}$, the hypergraph induced on $L_j$ is a permutation
hypergraph associated with $T-x$ and one parity class of bijections
$X\setminus\{x\}\to Y\setminus\{j\}$.

Put $c=r(x)$ and $J=Y\setminus\{c\}$. Choose a cyclic permutation $\sigma$
of $J$ such that
\begin{equation}\label{eq:sigma-condition}
 \sigma(r(b))\neq r(a).
\end{equation}
This is possible because $|J|=m-1\geq4$: a cyclic order can avoid one
prescribed directed adjacency. Set $r_c=r$. For every $j\in J$, choose
$r_j\in L_j$ satisfying
\begin{equation}\label{eq:rj}
 r_j(x)=j,\qquad r_j(a)=c,\qquad r_j(b)=\sigma(j).
\end{equation}
The three prescribed images are distinct. At least two positions and two
values remain, so the parity can be adjusted by interchanging two unused
images. Hence $r_j$ can always be chosen in $\Omega$.

By induction, each layer $L_j$ contains a loose hypertree $B_j$ spanning
$L_j\setminus\{r_j\}$. Initially that layer contributes the two components
$B_j$ and $R_j=\{r_j\}$. For $j\in J$, add
\begin{equation}\label{eq:Fj}
 F_j=\{r_j,r_j\tau,r_j\tau^2\}.
\end{equation}
The conditions in \eqref{eq:rj} imply
$
 r_j\tau\in L_c,$ and $ r_j\tau^2\in L_{\sigma(j)}.
$
Moreover, $r_j\tau\neq r_c$: equality would give
$j=r(b)$ and $\sigma(j)=r(a)$, contrary to
\eqref{eq:sigma-condition}. Also $r_j\tau^2\neq r_{\sigma(j)}$, since
$
 (r_j\tau^2)(a)=r_j(x)=j\neq c=r_{\sigma(j)}(a).
$
Thus $F_j$ meets the three components $R_j$, $B_c$, and
$B_{\sigma(j)}$.

Write the cycle of $\sigma$ as $(j_1\,j_2\,\cdots\,j_q)$. Add
$F_{j_1},F_{j_2},\ldots,F_{j_q}$ in this order. Before $F_{j_t}$ is
added, the growing component containing $B_c$ also contains
$B_{j_2},\ldots,B_{j_t}$ and $R_{j_1},\ldots,R_{j_{t-1}}$. The singleton
$R_{j_t}$ and the hypertree $B_{j_{t+1}}$, with indices read cyclically, are
still separate. Hence each new hyperedge joins three distinct
incidence-tree components and creates no incidence cycle. At the end every
vertex except $r_c=r$ belongs to one loose hypertree, while $r$ remains
isolated.
\end{proof}

\end{document}
